\documentclass[10pt,a4paper]{amsart}
\usepackage[margin=19mm]{geometry}
\usepackage{amsmath,amssymb,mathtools}
\usepackage[scr=rsfs,cal=euler]{mathalfa}
\usepackage{xfrac}
\usepackage[unicode,hidelinks,bookmarks=false]{hyperref}
\newcommand{\ie}{i.e.\ }
\newcommand{\cf}{cf.\ }
\newcommand{\resp}{respectively\ }
\newcommand{\Subsection}{\subsection}
\providecommand{\nobreakdash}{\nobreak\hbox{-}}
\setlength{\emergencystretch}{2em}
\multlinegap0pt
\usepackage{bm}
\usepackage{bbm}
\usepackage{dsfont}
\usepackage{xspace}
\usepackage{cancel}
\usepackage{mathtools}
\usepackage{amsthm}
\makeatletter
\@ifundefined{lemma}{\newtheorem{lemma}{Lemma}}{}
\@ifundefined{proposition}{\newtheorem{proposition}{Proposition}}{}
\makeatother
\title[Lumping of reaction networks: Generic and critical parameter]{Lumping of reaction networks: Generic and critical parameters}
\author{Justin Eilertsen, Valery G. Romanovski, Santiago Schnell, Sebastian Walcher}
\date{Source: arXiv:2606.28895v2 (CC BY 4.0)}
\begin{document}
\maketitle

\noindent\emph{Source and licence.} Lumping of reaction networks: Generic and critical parameters. Version arXiv:2606.28895v2 (CC BY 4.0), distributed under the Creative Commons Attribution 4.0 International licence, as recorded in the arXiv licence metadata for this exact version and shown on the abstract page https://arxiv.org/abs/2606.28895v2. Full rights notice: licenses/arxiv-2606.28895-CC-BY-4.0.txt.\par\smallskip

\noindent\textit{Editorial scope.} Counted unit: the Proposition whose source label is \texttt{colslem} in the article (source0.tex lines 1099--1101, source label \texttt{colslem}), delivered together with the article's own complete original proof of it (source0.tex lines 1102--1144, one \texttt{proof} environment of 43 lines). No mathematics is written, repaired or re-derived: statement and proof are the article's own text line for line. Required context retained uncounted: parode (lines 231--233); lalem (lines 800--806); rowech (lines 850--854); properlumpmat (lines 1094--1096). Every definition, display and label of those lines is retained unchanged. Omitted from this package and not redistributed: 1--230, 234--799, 807--849, 855--1093, 1097--1098, 1145--2414. Nothing inside a retained excerpt is omitted or altered: every retained body is delivered exactly as the article's lines are written, so no omission receipt of any kind accompanies this package. Citations: the retained excerpts carry no citation command at all, so no bibliography entry of the article is transcribed and no reference is redistributed. Reference adaptation: none was needed and none was made; every reference inside the delivered unit resolves inside it, so no unresolved reference or citation is produced. Printed slips kept literal: no printed word, symbol, hypothesis or number of the counted statement or of its proof was altered. Layout and class: the article's own class is \texttt{article} with packages including amssymb, verbatim, amsmath, amsfonts, amsthm, xcolor, array, enumerate, graphicx, mhchem, url, tikz-cd; the wrapper is the lane's pinned \texttt{amsart} editorial wrapper at 10pt on A4, which restates the theorem-like environments the retained text needs and adds the article's own macro definitions that the retained text needs (none), with extra wrapper packages (bm, bbm, dsfont, xspace, cancel, mathtools). The delivered numbering is the unit's own. Assets: no retained span contains a figure environment, a graphics-inclusion command, a PDF-page inclusion, a table or a TikZ picture; no image file is copied into this package, the package declares no asset dependency and no image file is redistributed with it. Licence: version arXiv:2606.28895v2 of this article is distributed under the Creative Commons Attribution 4.0 International licence, as recorded in the arXiv licence metadata for this exact version and shown on the abstract page https://arxiv.org/abs/2606.28895v2. A full rights notice is delivered at licenses/arxiv-2606.28895-CC-BY-4.0.txt. Characters: every retained excerpt is pure ASCII, so the delivered engine, XeLaTeX with its default Latin Modern text font, reports no missing character.
\begin{equation}\label{parode}
\dot x=F(x,k), \quad x\in \mathbb R^n,\, k\in \mathbb R_{+}^d,
\end{equation}

\begin{lemma}\label{lalem}
With $T\in\mathbb R^{e\times n}$ given, let $b_1,\ldots,b_{n-e}\in\mathbb R^n$ be a basis of $\ker T$, and let $B$ be the matrix with columns $b_i$. Then $k^*$ is a critical parameter for $T$ and \eqref{parode} if and only if 
\begin{equation}\label{linparcond}
T\,DF(x,k^*)\,B=0
\end{equation}
for all $x$.
\end{lemma} 

    \begin{equation}\label{rowech}
         QT=\begin{pmatrix}
        E_e & \widehat T
    \end{pmatrix},
    \end{equation}

\begin{equation}\label{properlumpmat}
 Tx=\begin{pmatrix} \sum_{j\in I_1} x_j \\ \vdots \\ \sum_{j\in I_r} x_j\end{pmatrix}.
\end{equation}

\begin{proposition}\label{colslem} Let system \eqref{parode} be given, and abbreviate $M=DF(x,k^*)$. With $T$ as in \eqref{properlumpmat}, and $1\leq p,q,\leq r$, denote by  $M_{pq}$  the $|I_p|\times |I_q|$ submatrix obtained from $M$ by deleting all rows with numbers not in $I_p$ and all columns with numbers not in $I_q$. \\
Then  $k^*$ is a critical parameter for $T$ and \eqref{parode} if, and only if, for each pair $(p,\, q)$ all column sums of $M_{pq}$ are equal.
\end{proposition}

\begin{proof} 
Let $B$ be such that its columns form a basis of $\ker T$. By Lemma \ref{lalem} it suffices to show that $M$ satisfies $T\cdot M\cdot B=0$ if and only if  all column sums of $M_{ij}$ are equal.\\
One may assume that 
\[
T=\begin{pmatrix} 1&\cdots&1&0&\cdots& 0& 0&\cdots\\
0&\cdots &0&1&\cdots&1&0&\cdots\\
\vdots&&&&&&&\ddots
\end{pmatrix},
\]
and that all rows with a single entry $1$ are gathered in the last columns. Letting
\[
M=\begin{pmatrix} M_{11}&\cdots&M_{1r}\\
          \vdots & \ddots&\vdots \\
               M_{r1}&\cdots & M_{rr}\end{pmatrix}
\]
according to the partitioning, one finds

\[
T\cdot M=\begin{pmatrix}{\rm cols}\,( M_{11})&\cdots&{\rm cols}\,(M_{1r})\\
          \vdots & \ddots&\vdots \\
               {\rm cols}\,(M_{r1})&\cdots & {\rm cols}\,(M_{rr})\end{pmatrix},
\]
where ${\rm cols}\,(M_{pq})$ denotes the row which has as entry $\#\ell$ the sum of the elements of column $\#\ell$ of $M_{pq}$.
Now (in a variant of Remark \ref{rowech}) the matrix $B$ built from basis elements of $\ker T$ can be chosen as
\[
B=\begin{pmatrix} B_1&0&\cdots& 0\\
                             0& \ddots &  & \vdots\\
                             \vdots& & \ddots &0\\
                              0& \cdots &  0& B_s
\end{pmatrix},
\]
with each
\[
B_j=\begin{pmatrix} 1&0&0&\cdots& 0\\
                               -1&1&0&\cdots &0\\
0&-1& 1 && 0  \\
\vdots& & \ddots &\ddots & \vdots\\
0&\cdots &\cdots & -1& 1 \\
0&\cdots &&0 &-1
\end{pmatrix}
\]
of appropriate size, corresponding to a row of $T$ with more than one entry $1$. Now multiplication of $TM$ by $B$ shows, for each index pair $(p,\,q)$, that all entries of ${\rm cols}\,(M_{pq})$ are equal.
\end{proof}

\end{document}
